% GENERATED FILE. Edit canonical lessons, then run npm run generate:books.
% canonical-source-hash: fnv1a64:594e6bbece36849a
% canonical-lessons: GE-C253-dann, GE-C253-schliesslich, GE-C253-lesung, GE-C253-review

\chapter{Zuerst, Dann, Schließlich — A Story in Order}
\label{ch:ge-a-story-in-order}
\hlchaptermodality{253}

\begin{chapteropening}
\textbf{By the end of this chapter:} \emph{I can tell a short connected story in order with zuerst, dann, danach and schließlich, and read one of about a hundred words.}

Complete GE-C253-review: retell the reading's journey aloud in order, with the verb second after every step word.
\end{chapteropening}

\section[zuerst, dann, danach]{dann — and the verb comes next}

\label{lesson:GE-C253-dann}



\begin{quote}
Say \textquotedblleft{}first\textquotedblright{} and \textquotedblleft{}afterwards\textquotedblright{}. (\emph{Zuerst. Danach.})

A story needs one more step word between those two: \textbf{dann}, \textquotedblleft{}then\textquotedblright{}.
\end{quote}

\begin{grammarlens}[title={Grammar Lens: the step word, then the verb}]
In a German statement the verb is the second piece. When a step word takes the first place, the verb comes straight after it, and the subject moves behind the verb:

\begin{quote}
\textbf{Zuerst} \emph{kaufte ich einen Kaffee.} — First I bought a coffee. \textbf{Dann} \emph{las ich die Nachricht.} — Then I read the message. \textbf{Danach} \emph{ging ich.} — Afterwards I went.
\end{quote}

English keeps \textquotedblleft{}I\textquotedblright{} in front: \textquotedblleft{}then I read\textquotedblright{}. German says \emph{dann las ich}, not \emph{dann ich las}. The verb holds its place, and the subject gives way.

\emph{Und} does not take a place of its own. After \emph{und dann}, the verb still follows \emph{dann}: \emph{… und dann kam der Zug.}

\emph{Dann} and \emph{danach} overlap. \emph{Dann} is the next step, and the shortest way on. \emph{Danach} leans on the time: after that.
\end{grammarlens}

\begin{cousinweb}
\textbf{dann} $\leftarrow$ Old High German \emph{danne} $\leftarrow$ Proto-Germanic \emph{*þan}, \textquotedblleft{}at that time\textquotedblright{} (EtymDB 2.1), which EtymDB sets beside English \emph{then}. In chapter forty-eight you met its twin, \textbf{denn}, \textquotedblleft{}because\textquotedblright{}: the two were once one word, and German gave \emph{dann} the when and \emph{denn} the why.
\end{cousinweb}

\subsection*{Your turn}
Say these aloud:

\begin{itemize}
\raggedright
  \item \emph{Zuerst kaufte ich einen Kaffee.}
  \item \emph{Dann las ich die Nachricht.}
  \item \emph{Danach ging ich.}
  \item the three in a row, without stopping
  \item \emph{… und dann kam der Zug.} The verb after \emph{dann}, not after \emph{und}
\end{itemize}

\begin{tcolorbox}[breakable,colback=teal!4,colframe=teal!35!black,title={Before you move on}]
\textquotedblleft{}Then I went\textquotedblright{}? (\emph{\textbf{Dann ging ich.}}) What comes straight after \emph{dann}? (\textbf{The verb.}) Does \emph{und} change that? (\textbf{No:} \emph{\textbf{und dann kam er.}}) Which old word is \emph{dann}'s twin? (\emph{\textbf{Denn.}})
\end{tcolorbox}
\section[schließlich]{schließlich — the step that closes the story}

\label{lesson:GE-C253-schliesslich}



\begin{quote}
Say \textquotedblleft{}then I went\textquotedblright{}. (\emph{Dann ging ich.}) Now \textquotedblleft{}first I bought a coffee\textquotedblright{}. (\emph{Zuerst kaufte ich einen Kaffee.})
\end{quote}

\subsection*{You'll want to know: schließlich}
\textbf{schließlich} — \textquotedblleft{}in the end, finally\textquotedblright{}: \emph{schlees-likh}.

\begin{quote}
\emph{Schließlich kam der Zug.} — In the end the train came.
\end{quote}

You already have \textbf{endlich}, \textquotedblleft{}at last\textquotedblright{}. The two are close, and the difference is the waiting. \emph{Endlich} sighs: at last, after too long. \emph{Schließlich} counts: the last step, in order.

\emph{Schließlich} has a second life outside stories, \textquotedblleft{}after all\textquotedblright{}: \emph{Er ist schließlich mein Bruder}, he is my brother, after all.

\begin{grammarlens}[title={Grammar Lens: the frame of a told story}]
Four step words, and every one takes the first place, so the verb follows it at once:

\begin{quote}
\textbf{Zuerst} \emph{wartete ich.} \textbf{Dann} \emph{kam der Zug.} \textbf{Danach} \emph{fuhr ich zwei Stunden.} \textbf{Schließlich} \emph{kam ich in Berlin an.}
\end{quote}

Around those steps sits the scene: \emph{Es war Samstag. Es regnete.} A told story is the scene, the steps in order, and the end. You do not need all four step words every time: \emph{zuerst … dann … schließlich} is already a story.
\end{grammarlens}

\begin{cousinweb}
\textbf{schließlich} is \textbf{schließen}, \textquotedblleft{}to close\textquotedblright{}, + \emph{-lich}. \emph{Schließen} goes back through Middle High German \emph{sliezen} to Old High German \emph{sliozan}, \textquotedblleft{}to close, to conclude\textquotedblright{} (EtymDB 2.1). EtymDB records \emph{schließlich} without taking it apart, so the split into verb and ending is given here as the plain reading of the word: the closing one.
\end{cousinweb}

\subsection*{Your turn}
Say these aloud:

\begin{itemize}
\raggedright
  \item \emph{schließlich}
  \item \emph{Schließlich kam der Zug.}
  \item the four steps: \emph{Zuerst wartete ich. Dann kam der Zug. Danach fuhr ich zwei Stunden. Schließlich kam ich in Berlin an.}
  \item \emph{endlich}, at last, after waiting; \emph{schließlich}, the last step
  \item three steps of your own morning, with \emph{zuerst}, \emph{dann}, \emph{schließlich}
\end{itemize}

\begin{tcolorbox}[breakable,colback=teal!4,colframe=teal!35!black,title={Before you move on}]
\textquotedblleft{}In the end I arrived in Berlin\textquotedblright{}? (\emph{\textbf{Schließlich kam ich in Berlin an.}}) Which word sighs after waiting, and which counts the last step? (\emph{\textbf{Endlich}} sighs; \emph{\textbf{schließlich}} counts.) Which verb is inside \emph{schließlich}? (\emph{\textbf{Schließen}}, to close.)
\end{tcolorbox}
\section[(reading)]{(reading) — a journey, told in order}

\label{lesson:GE-C253-lesung}



\begin{quote}
Chapter fifty-two gave you forty-eight words in a row, in the present. This passage is twice as long, and it happened last Saturday.

Every word in it has come up before this lesson.
\end{quote}

\begin{tcolorbox}[breakable,skin=enhanced,colback=orange!6,colframe=orange!45!black,boxrule=0.5pt,arc=1mm,left=6pt,right=6pt,top=4pt,bottom=4pt,fonttitle=\bfseries,title={Read this}]
\begin{quote}
Es war Samstag, und es regnete. Ich wartete am Bahnhof, aber der Zug hatte Verspätung. Zuerst kaufte ich einen Kaffee. Dann las ich die Nachrichten auf meinem Handy. Als der Zug endlich kam, war ich sehr müde.
\end{quote}

\begin{quote}
Ich fuhr zwei Stunden. Danach kam ich in Berlin an. Dort wartete meine Großmutter. Sie lachte, als sie mich sah. Wir gingen zusammen in ein Restaurant. Ich hatte Hunger und nahm einen Fisch, und meine Großmutter nahm eine Suppe. Sie erzählte mir eine Geschichte, und wir lachten. Schließlich war es Abend. Ich war müde, aber sehr glücklich.
\end{quote}

Read it once without stopping. Do not translate as you go.

Then read it a second time and listen for the step words that put the day in order.
\end{tcolorbox}

\begin{culture}
Four step words carry the day: \emph{zuerst}, \emph{dann}, \emph{danach}, \emph{schließlich}. After each one the verb comes next: \emph{zuerst kaufte ich}, \emph{dann las ich}, \emph{danach kam ich}, \emph{schließlich war es}. \emph{Dort} does the same job in \emph{Dort wartete meine Großmutter}.

Every verb is a Präteritum. Some take \emph{-te}: \emph{regnete}, \emph{wartete}, \emph{kaufte}, \emph{lachte}, \emph{erzählte}. Some change their vowel: \emph{las}, \emph{kam}, \emph{fuhr}, \emph{sah}, \emph{gingen}, \emph{nahm}. And \emph{hatte} and \emph{war} you have said since chapters twenty-five and twenty-eight.

Two lines turn on \textbf{als}. \emph{Als der Zug endlich kam} is one moment that ended the waiting; \emph{als sie mich sah} is the moment the grandmother laughed. The first stands in front, so verb meets verb at the comma: \emph{kam, war}. The second follows its main clause: \emph{Sie lachte, als sie mich sah}.
\end{culture}

\subsection*{Your turn}
\begin{itemize}
\raggedright
  \item \emph{Say it:} the passage aloud, once, without stopping
\end{itemize}

\begin{tcolorbox}[breakable,colback=teal!4,colframe=teal!35!black,title={Before you move on}]
How many words in that passage were new? (\textbf{None.})

Why does \emph{Dann las ich} put \emph{las} before \emph{ich}? (\emph{Dann} takes the \textbf{first place}, and the verb comes \textbf{second}.)

Which word marks the one moment the train came? (\emph{\textbf{Als.}})
\end{tcolorbox}
\section[(review)]{Review — the journey, retold}

\label{lesson:GE-C253-review}



\begin{quote}
Say \textquotedblleft{}then I read the message\textquotedblright{}. (\emph{Dann las ich die Nachricht.}) And \textquotedblleft{}in the end it was evening\textquotedblright{}. (\emph{Schließlich war es Abend.})
\end{quote}

\begin{grammarlens}[title={What you've built}]
Two chapters, and a German story is yours to tell. The scene in the Präteritum: \emph{Es war Samstag, und es regnete.} The steps, each with its verb straight after it: \emph{zuerst}, \emph{dann}, \emph{danach}, \emph{schließlich}. One moment marked with \emph{als}. And a last line about how it felt.

Now tell the passage back, without reading it. It does not have to be word for word. It has to be in order, and every step word has to pull its verb in behind it.
\end{grammarlens}

\subsection*{Your turn}
Say these aloud:

\begin{itemize}
\raggedright
  \item the scene: \emph{Es war Samstag, und es regnete.}
  \item the waiting: \emph{Ich wartete, aber der Zug hatte Verspätung.}
  \item \emph{Zuerst kaufte ich einen Kaffee. Dann las ich die Nachrichten.}
  \item \emph{Als der Zug endlich kam, war ich sehr müde.}
  \item \emph{Danach kam ich in Berlin an.} Then one more line of your own about the grandmother
  \item the end: \emph{Schließlich war es Abend.}
\end{itemize}

\begin{tcolorbox}[breakable,colback=teal!4,colframe=teal!35!black,title={Before you move on}]
The four step words, in order? (\emph{\textbf{Zuerst, dann, danach, schließlich.}}) After each one, what comes next? (\textbf{The verb.}) Which one closes the story? (\emph{\textbf{Schließlich.}})
\end{tcolorbox}
